\documentclass{article}
\usepackage{amsmath, amssymb, amsthm}
\usepackage{hyperref}
\usepackage{geometry}
\geometry{margin=1in}

\title{Erd\H{o}s Problem \#1113}
\date{Last updated: 2025-12-28}
\author{Source: erdosproblems.com}

\begin{document}
\maketitle

\noindent\textbf{Status:} OPEN \quad \textbf{No prize}

\noindent\textbf{Tags:} number theory, covering systems

\medskip
\noindent\textbf{Problem Statement:}

A positive odd integer $m$ such that none of $2^km+1$ are prime for $k\geq 0$ is called a Sierpinski number. We say that a set of primes $P$ is a covering set for $m$ if every $2^km+1$ is divisible by some $p\in P$.

Are there Sierpinski numbers with no finite covering set of primes?

\bigskip
\noindent\textbf{Remarks:}

Sierpinski \cite{Si60} proved that there are infinitely many Sierpinski numbers, using covering systems to construct suitable covering sets for any $m$ satisfying a certain congruence. This establishes that there is a positive density set of such $m$.

The smallest Sierpinski number is believed to be $78557$, which was found by Selfridge.

Erd\H{o}s and Graham \cite{ErGr80} asked whether there are Sierpinski numbers for which a covering system is not 'responsible', for which the best interpretation seems to be the above question. This is formulated precisely in problem F13 of Guy's collection \cite{Gu04}. Erd\H{o}s and Graham thought the answer is yes (in that there are such Sierpinski numbers), since otherwise this would imply there are infinitely many Fermat primes.

There is now further evidence with a concrete example: an argument of Izotov \cite{Iz95}, given in more detail by Filaseta, Finch, and Kozek \cite{FFK08}, suggests that $m=734110615000775^4$ is a Sierpinski number without a covering set. (Izotov proved that this $m$ is indeed a Sierpinski number.)

Filaseta, Finch, and Kozek \cite{FFK08} give a revised conjecture, suggesting that every Sierpinski number is either a perfect power or else has a finite covering set of primes. They also prove that for every $l\geq 1$ there is an $m$ such that $2^km^i+1$ is composite for all $1\leq i\leq l$ and $k\geq 0$.

See also [203], and [276] for another problem in which the question is whether covering systems are always responsible.

\bigskip
\noindent\textbf{References:}

[ErGr80] Erd\H{o}s, P. and Graham, R., Old and new problems and results in combinatorial number theory. Monographies de L'Enseignement Mathematique (1980).
 
 [FFK08] Filaseta, Michael and Finch, Carrie and Kozek, Mark, On powers associated with Sierpi\'nski numbers, {R}iesel
numbers and {P}olignac's conjecture. J. Number Theory (2008), 1916--1940.
 
 [Gu04] Guy, Richard K., Unsolved problems in number theory. (2004), xviii+437.
 
 [Iz95] Izotov, Anatoly S., A note on Sierpi\'nski numbers. Fibonacci Quart. (1995), 206--207.
 
 [Si60] Sierpi\'nski, W., Sur un probl\`eme concernant les nombres {$k\cdot 2\sp{n}+1$}. Elem. Math. (1960), 73--74.

\bigskip
\noindent\small{Source: \url{https://www.erdosproblems.com/1113}}
\end{document}
